% Flooding advantage and growth-rate trade-off --- plan \S4.12 (Phase 3).
% Decision: stays in Ch.4 (chapter's population-level climax); Ch.6 takes
% the evolutionary discussion.
% Sources: new_notes3 Part E \S22--\S25; new_notes Part II \S13;
% new_notes2 \S9. All displayed numerical values re-verified 2026-08-07.

\section{The flooding advantage and the growth-rate trade-off}
\label{sec:flooding}

Does bursting help a pathogen get established? The question belongs to the
stochastic early phase, where censuses are small and extinction is a real
threat, and it can now be answered exactly: both the bursting and the
budding version of the early infection are branching processes whose
offspring distributions are available in closed form.

\subsection{Offspring of one infected cell}

In the branching approximation each released particle independently either
infects a target cell (rate $\gamma T$) or is cleared (rate $c$), so it
infects with probability
\begin{equation}
q=\frac{\gamma T}{\gamma T+c}.
\end{equation}
The offspring of one infected cell is the burst size $\Kb$, thinned by this
success probability. The (unconditional) burst-size \pgf{} is, by
Theorem~\ref{thm:burst},
$G_{\Kb}(z)=b+(\delta/\beta)\,z/(a-z)$ --- including the mass $b$ of cells
that never burst --- and thinning replaces $z$ by $y:=1-q+qz$:
\begin{equation}
G_{\rm off}(z)=b+\frac{\delta}{\beta}\,\frac{y}{a-y},
\qquad y=1-q+qz.
\label{eq:Goff}
\end{equation}
One checks $G_{\rm off}(1)=b+(\delta/\beta)/(a-1)=b+B=1$. The mean and
variance of the offspring number are
\begin{equation}
m=G_{\rm off}'(1)=q\,V_\infty,
\qquad
\mathrm{Var}_{\rm burst}
=\frac{2q^2V_\infty}{a-1}+qV_\infty-q^2V_\infty^2,
\label{eq:offmoments}
\end{equation}
the variance following from $\mathrm{Var}(\Kb)$ via the limit of
\eqref{eq:EW2} and the standard composition formula for random sums. When
$\gamma T\ll c$, $q\approx\gamma T/c$ and $m\approx R_0^{\mathrm{ODE}}$ of
\eqref{eq:R0ODE}.

\subsection{Extinction probability in closed form}

The extinction probability of a branching process is the smallest root of
$G_{\rm off}(z)=z$ in $[0,1]$. With $y=1-q+qz$, the fixed-point equation is
quadratic, $z=1$ is always a root, and the other root gives, for $m>1$,
\begin{equation}
z_{\rm ext}^{\rm burst}
=\frac{a(1-q+qb)-1+q}{q}
=\frac{a-1}{q}+1-a(1-b).
\label{eq:zext}
\end{equation}

\subsection{Comparison with matched-mean budding: the main theorem}

The comparator is the classical picture: a cell with $\mathrm{Exp}(d_{\Icell})$
lifetime releasing at constant rate $p$. The number it releases is a Poisson
process stopped at an independent exponential time, hence geometric on
$\{0,1,2,\ldots\}$ with mean $p/d_{\Icell}$; matched to the bursting cell's
total output, $p/d_{\Icell}=V_\infty$. Thinning a geometric law leaves a
geometric law, so the budding offspring number is geometric with the same
mean $m=qV_\infty$, and its extinction probability is
\begin{equation}
z_{\rm ext}^{\rm bud}=\frac{1}{m}\qquad(m>1).
\end{equation}

\begin{theorem}[The flooding advantage]
\label{thm:flood}
Write $L:=a(1-b)=a-\mu/\beta$, so that $V_\infty=L/(a-1)$. For $m>1$,
\begin{equation}
z_{\rm ext}^{\rm burst}-z_{\rm ext}^{\rm bud}
=\bigl(L-1\bigr)\Bigl(\frac{1}{m}-1\Bigr).
\label{eq:flood}
\end{equation}
Hence bursting strictly lowers the establishment extinction probability,
relative to matched-mean budding, if and only if $L>1$, i.e.
\begin{equation}
\delta(\beta+\mu)>\beta\mu
\qquad\Longleftrightarrow\qquad
\boxed{\;\frac{1}{\delta}<\frac{1}{\beta}+\frac{1}{\mu}\;}
\end{equation}
(with the convention that the condition holds for all $\delta>0$ when
$\mu=0$). When $L=1$ the two extinction probabilities are identical for
every $q$.
\end{theorem}

\begin{proof}[Proof sketch]
The identity \eqref{eq:flood} is algebra: substitute \eqref{eq:zext},
$m=qV_\infty=qL/(a-1)$ and simplify. For the criterion, $L>1$ is
$a(1-b)>1$, i.e.\ $a>1+\mu/\beta$; inserting the definition of $a$ and
squaring (both sides are positive, or the right-hand side is not) yields
$(\beta+\mu+\delta)^2-4\beta\mu>4\beta^2(1+\mu/\beta-\eta)^2$, which
reduces to $\delta(\beta+\mu)>\beta\mu$.
\end{proof}

Numerically, the identity \eqref{eq:flood} is exact in every case we
checked:
\begin{center}
\renewcommand{\arraystretch}{1.2}
\begin{tabular}{@{}llccc@{}}
\toprule
$(\beta,\mu,\delta)$ & $q$ & $z_{\rm ext}^{\rm burst}$ & $z_{\rm ext}^{\rm bud}$ & verdict \\
\midrule
$(1,0,0.1)$ & $0.2$ & $0.400$ & $0.455$ & bursting wins ($L=1.10$) \\
$(1,0.9,0.1)$ & $0.9$ & $0.935$ & $0.844$ & budding wins ($L=0.42$) \\
$(1,0.5,1/3)$ & $0.6$ & $0.833$ & $0.833$ & tie ($L=1$) \\
$(1,0.5,1/3)$ & $0.9$ & $0.556$ & $0.556$ & tie ($L=1$) \\
\bottomrule
\end{tabular}
\end{center}
For $m\le1$ both processes die out almost surely and the comparison is
empty; \eqref{eq:flood} changes sign there, an artefact of both formulae
running outside their domain $[0,1)$, not a reversal of the biology.

The interpretation is that the flooding advantage is real but bought
against a risk. Bursting concentrates output into one batch --- lowering
the offspring variance relative to the geometric output of a budding cell
--- but it carries the all-or-nothing hazard that the intracellular
population dies out first, with probability $b$, and the cell releases
nothing. The condition says the advantage survives exactly when the burst
timescale $1/\delta$ is shorter than the harmonic-type combination
$1/(1/\beta+1/\mu)$ of the birth and death timescales. When $\mu=0$ the
hazard of internal extinction vanishes and bursting always wins.

\subsection{The boundary $L=1$: a coincidence of laws}

The boundary case is more than a tie of extinction probabilities.

\begin{proposition}
\label{prop:L1}
The following are equivalent:
\begin{equation}
L=1
\;\Longleftrightarrow\;
a=1+\frac{\mu}{\beta}
\;\Longleftrightarrow\;
\delta(\beta+\mu)=\beta\mu
\;\Longleftrightarrow\;
\frac{1}{\delta}=\frac{1}{\beta}+\frac{1}{\mu}.
\end{equation}
When they hold, the bursting offspring law \eqref{eq:Goff} is
\emph{exactly geometric} with mean $m$ --- not merely equal in extinction
probability to the budding comparator. In particular the offspring
variances coincide, and
\begin{equation}
\mathrm{Var}_{\rm bud}-\mathrm{Var}_{\rm burst}
=\frac{2q^2V_\infty(L-1)}{a-1},
\label{eq:varorder}
\end{equation}
so that at matched mean $m$ the bursting offspring variance is smaller
than the budding one when $L>1$, equal when $L=1$, and larger when $L<1$.
\end{proposition}

\begin{proof}
$L=1$ is $a(1-b)=1$, i.e.\ $a-1=ab=\mu/\beta$, the chain of equivalences
above. Under it, $b=\delta/\beta$ (by \eqref{eq:AB}), and \eqref{eq:Goff}
collapses:
$$
G_{\rm off}(z)=b\Bigl(1+\frac{y}{a-y}\Bigr)=\frac{ab}{a-y}
=\frac{a-1}{a-y}
=\frac{(a-1)/a}{1-y/a},
$$
using $ab=a-1$. With $y=1-q+qz$,
$$
1-\frac{y}{a}=\frac{a-1+q}{a}\Bigl(1-\frac{q}{a-1+q}\,z\Bigr),
$$
so
$$
G_{\rm off}(z)=\frac{1-r}{1-rz},
\qquad
r=\frac{q}{a-1+q}=\frac{m}{1+m},
\qquad m=\frac{q}{a-1}=qV_\infty,
$$
which is the \pgf{} of a geometric law on $\{0,1,2,\ldots\}$ with mean $m$
--- the same law as the matched budding comparator. The variance identity
\eqref{eq:varorder} is the difference between the geometric variance
$m(1+m)$ and \eqref{eq:offmoments}, with $V_\infty=L/(a-1)$.
\end{proof}

At the boundary, then, bursting and budding are the same branching process
in different clothes. Away from it, \eqref{eq:varorder} quantifies the
direction in which they differ.

\subsection{The deterministic growth-rate trade-off}
\label{sec:r-tradeoff}

Establishment is only half of the story; the other half is how fast an
established infection grows. Match the two models a second time, now at
level $R_0$ and at mean productive lifetime: take the budding cell with
$d_{\Icell}=1/\mean{T_{\rm prod}}$ and $p=V_\infty d_{\Icell}$ --- which
is exactly $p_{\rm eff}(0)$, so this is the fairest possible classical
match --- and choose $\gamma T$ (or $c$) so that both models share
$R_0=2$. Then solve for the deterministic growth rates: $r_{\rm bud}$ from
the quadratic $(r+d_{\Icell})(r+c)=\gamma Tp$, and $r_{\rm burst}$ from
the characteristic equation \eqref{eq:char}. In every regime we checked
--- $L>1$, $L=1$, and $L<1$ --- the ordering is the same:
\begin{equation}
r_{\rm bud}>r_{\rm burst}.
\end{equation}
With $c=1$ and $R_0=2$:
\begin{center}
\renewcommand{\arraystretch}{1.2}
\begin{tabular}{@{}lcccc@{}}
\toprule
$(\beta,\mu,\delta)$ & $L$ & $d_{\Icell}$ & $r_{\rm bud}$ & $r_{\rm burst}$ \\
\midrule
$(1,0,0.1)$ & $1.10$ & $0.417$ & $0.250$ & $0.180$ \\
$(1,0.5,1/3)$ & $1.00$ & $0.910$ & $0.395$ & $0.294$ \\
$(1,0.9,0.1)$ & $0.42$ & $0.701$ & $0.343$ & $0.198$ \\
\bottomrule
\end{tabular}
\end{center}
The mechanism is generation time: the budding cell starts releasing
immediately, while the bursting cell must first accumulate load, so
budding's generation times are shorter and its exponential invasion is
faster --- the same delay that appeared as the eclipse-like lag of
Figures~\ref{fig:overlay-V} and \ref{fig:growth-phase}.

The two comparisons together give the biological headline. At fixed total
output per cell, bursting can lower the probability of early extinction
(when $L>1$) while simultaneously slowing deterministic invasion speed:
$R_0$ stays fixed, and $z_{\rm ext}$ and $r$ move in opposite directions.
A pathogen tuning its burst rate $\delta$ faces a genuine trade-off. For
$\mu>0$, budding beats bursting on extinction until $\delta$ crosses the
threshold $\delta_*=\beta\mu/(\beta+\mu)$ --- the $L=1$ boundary --- beyond
which bursting takes over; for $\mu=0$ bursting wins on extinction at every
$\delta$, at the continuing cost of a slower $r$.

\begin{remark}[Scope of the comparison]
Both results compare branching structures at fixed offspring mean. They say
nothing about nonlinear effects --- immune saturation, local depletion of
targets, Allee-type cooperation --- where other advantages of flooding
may live; any full account should name and separate the two mechanisms.
The evolutionary framing, and where this result sits relative to the
bursting--budding spectrum of Chapter~6, is taken up there.
\end{remark}
